\chapter{Taking the model to data}\label{ch:empirical}

\begin{chapterplan}
Which comparisons of prices, order counts and outcomes can a researcher compute from public data, and which sign does
the model predict for each? Every result concerns one target in one public state: the prior, the slot and the
incumbent law stay fixed (Appendix A.8 of the paper).
\end{chapterplan}

\section{Dating and full two-sided trade}\label{sec:empirical-dating}

The interval runs from the public sale news to the challenger's decision (substantive diligence, not the first
public offer); price and strength come before it. The order \(X\) is the sign of the net imbalance at the first
prespecified opportunity, such as a scheduled opening auction, with a zero when no order arrives.

Under \emph{full two-sided trade} (\(\sigma_H=\delta_1\), \(\sigma_L=\delta_{-1}\)), the order frequencies are
\(\bar f(1)=b+\pi\lambda\), \(\bar f(-1)=b+(1-\pi)\lambda\), the volume is \(v=2b+\lambda\), and
\(\E[X]=\lambda(2\pi-1)\). Bayes' rule gives the \emph{balance identity}
\begin{equation}\label{eq:empirical-balance}
\bar f(1)\,(M-\pi)=\bar f(-1)\,(\pi-m)=\pi(1-\pi)\lambda .
\end{equation}
A Bayesian belief has no drift, so the expected rise after a purchase equals the expected fall after a sale. The
impacts are \(I_+=P(1)-P(0)\) and \(I_-=P(0)-P(-1)\), and \(w_F(\mu)=t_L-t_0+\DT\mu\).

\section{The price gap and the weighted test (Proposition A.10)}\label{sec:empirical-spread}

{\small
\begin{paperresult}{Proposition A.10 (The price gap and the weighted test)}
Standing hypotheses of Appendix A.8: \(F\) is any continuous law, the classes satisfy the class ranking, \(\pi\in(0,1)\) is arbitrary, \(\DT>0\) and \(\gamma=g_H-g_L>0\).
\begin{enumerate}[label=(\roman*),leftmargin=2em,nosep]
\item \emph{Every trading pattern.} In every equilibrium,
\(
\bar f(1)I_+-\bar f(-1)I_-=\E[P(X)]-P(0)=e(0)\,\DT\,[\pi-\mu(0)]+\sum_x\bar f(x)\,[e(x)-e(0)]\,w_F(\mu(x)).
\)
\item \emph{The weighted test.} Under full two-sided trade,
\(
\bar f(1)I_+-\bar f(-1)I_-=\bar f(1)\,[e(1)-e(0)]\,w_F(M)-\bar f(-1)\,[e(0)-e(-1)]\,w_F(m).
\)
If also \(\pi<\tau\), then \(I_+-\dfrac{\bar f(-1)}{\bar f(1)}\,I_-=[e(1)-\rho]\,w_F(M)\), which is positive exactly when a purchase recruits, that is, exactly when \(e(1)>\rho\), under both challenger rules.
\item \emph{Signs.} Let full two-sided trade hold, and let an indifferent high-cost challenger prepare. Entry after the orders \((-1,0,1)\) and the weighted impact difference are:
\begin{center}\normalfont\small
\begin{tabular}{@{}lll@{}}
\toprule
Threshold & Entry after \((-1,0,1)\) & \(\bar f(1)I_+-\bar f(-1)I_-\)\\
\midrule
\(\tau\le m\) & \((1,1,1)\) & \(0\)\\
\(m<\tau\le\pi\) & \((\rho,1,1)\) & \(-(1-\rho)\,\bar f(-1)\,w_F(m)<0\)\\
\(\pi<\tau\le M\) & \((\rho,\rho,1)\) & \((1-\rho)\,\bar f(1)\,w_F(M)>0\)\\
\(M<\tau\) & \((\rho,\rho,\rho)\) & \(0\)\\
\bottomrule
\end{tabular}
\end{center}
When an indifferent high-cost challenger may prepare with any probability, the table holds for \(\tau\notin\{m,\pi,M\}\). At \(\tau=M\) the weighted difference is \([e(1)-\rho]\,\bar f(1)\,w_F(M)\ge0\), positive exactly when \(e(1)>\rho\). At \(\tau=m\) it is \(-[1-e(-1)]\,\bar f(-1)\,w_F(m)\le0\), negative exactly when \(e(-1)<1\). At \(\tau=\pi\), full two-sided trade is an equilibrium with one entry \(e(0)\in[\rho,1]\) exactly when it is an equilibrium with every such entry. As \(e(0)\) runs over \([\rho,1]\), the weighted difference takes every value from \(-(1-\rho)\,\bar f(-1)\,w_F(m)\) to \((1-\rho)\,\bar f(1)\,w_F(M)\).
\item \emph{The price gap.} Under full two-sided trade with \(\pi<\tau\),
\(
P(1)-P(-1)=\rho\DT\,(M-m)+[e(1)-\rho]\,w_F(M),\ \rho\DT\,(M-m)=\frac{v}{\bar f(1)}\,I_-,
\)
a class component and an entry component. If the trading margin \(\rho\kappa(\pi)\DT-k\) is positive, with \(\kappa(\pi)=\min\{m,1-M\}\), every equilibrium has full two-sided trade, under both investor rules, whether or not the challenger sees the price. If also \(\pi<\tau\) and the challenger does not see the price, its belief stays at \(\pi\), entry is \(\rho\) after every order, the price gap is the class component \(\rho\DT\,(M-m)\), and the weighted impact difference is zero.
\item \emph{The unweighted comparison.} Under full two-sided trade with \(\pi<\tau\),
\(
I_+-I_-=\rho\DT\,\frac{\pi(1-\pi)\lambda^2(1-2\pi)}{\bar f(1)\,\bar f(-1)}+[e(1)-\rho]\,w_F(M).
\)
If \(\pi\ge1/2\), \(\rho\le1/2\) and an indifferent high-cost challenger prepares, then under full two-sided trade \(I_+>I_-\) exactly when \(\pi<\tau\le M\), that is, exactly when a purchase raises entry above its level after no order.
\item \emph{Strength.} Let \(F_r\) be a path along which \(\DT\) rises strictly. Hold the prior and the slot fixed, and write \(w_r=w_{F_r}\) and \(\chi_r=\chi_{F_r}\).
\begin{enumerate}[label=(\alph*),leftmargin=2em,nosep]
\item For every belief \(\mu\), \(w_r(\mu)=pF_r(p)+\int_p^\infty\bigl[1-(1-\mu)G_L(y)-\mu G_H(y)\bigr]\bigl[1-F_r(y)\bigr]\,dy\). From \(r\) to \(r'>r\), the first term falls by \(p\,[F_r(p)-F_{r'}(p)]\ge0\), and the second rises by \(\int_p^\infty[1-(1-\mu)G_L(y)-\mu G_H(y)]\,[F_r(y)-F_{r'}(y)]\,dy\ge\mu\,[\DT(r')-\DT(r)]\). The entry gain \(w_r(\mu)\) rises exactly when the rise of the second term exceeds the fall of the first.
\item If \(F_r(p)\) is the same at every strength of the path, \(w_r(\mu)\) rises strictly in \(r\) for every \(\mu\in(0,1]\). If \(F_r\) is uniform on \([0,r]\) with \(r>p\), then \(\frac{d}{dr}\,w_r(\mu)=\frac{\E_\mu[\min\{\theta,r\}^2]-3p^2}{2r^2}\), where \(\E_\mu\) takes \(\theta\) from \((1-\mu)G_L+\mu G_H\). So \(w_r(\mu)\) rises strictly in \(r\) for every \(\mu\) when \(p<\min\{\underline\theta,\underline r\}/\sqrt3\), where \(\underline\theta\) is the lowest challenger value and \(\underline r\) the lowest strength on the path.
\item Let an indifferent high-cost challenger prepare, and let full two-sided trade hold at every strength of an interval on which \(\pi<\tau(r)\) and \(w_r(M)\) rises strictly. Where \(\chi_r=1\), both components of the price gap, both impacts, the price gap and the weighted impact difference rise strictly with \(r\). Where \(\chi_r=0\), the entry component and the weighted impact difference are zero, and the class component, both impacts and the price gap rise strictly with \(r\).
\item The indicator \(\chi_r\) is nonincreasing in \(r\). Let also \(F_r(y)\) be continuous in \(r\) at every \(y\). Then the strengths with \(\chi_r=1\) form a closed subset of the path. Let \(r_C\) be a strength with \(\chi_{r_C}=1\) and \(\chi_r=0\) for every \(r>r_C\), and let full two-sided trade and \(\pi<\tau(r)\) hold on a neighborhood of \(r_C\). When an indifferent high-cost challenger prepares, \(P(1)\), \(I_+\) and the price gap are continuous from the left at \(r_C\), and their right limits lie below their values at \(r_C\) by \((1-\rho)\,w_{r_C}(M)>0\). The weighted impact difference falls by \((1-\rho)\,\bar f(1)\,w_{r_C}(M)\) in the same way. \(P(-1)\), \(P(0)\) and \(I_-\) are continuous at \(r_C\). When an indifferent investor may mix over its best orders and an indifferent high-cost challenger may prepare with any probability, entry after a purchase at \(r_C\) can take any value in \([\rho,1]\), and the price gap at \(r_C\) any value between its two one-sided limits.
\end{enumerate}
\end{enumerate}
\end{paperresult}}

{\small\begin{derivation}[proof]
(i) \(\E[P(X)]-P(0)=\sum_x\bar f(x)[P(x)-P(0)]=\bar f(1)I_+-\bar f(-1)I_-\). With \(P(x)=t_0+e(x)w_F(\mu(x))\) and
\(\sum_x\bar f(x)w_F(\mu(x))=w_F(\pi)\) (\(w_F\) is affine, the mean posterior is \(\pi\)), add and subtract
\(e(0)w_F(\pi)\). (ii)--(iii) Full trade gives \(\mu(0)=\pi\); for \(\pi<\tau\), \(e(-1)=e(0)=\rho\). The class parts
\(\rho\DT(M-\pi)\) of \(I_+\) and \(\rho\DT(\pi-m)\) of \(I_-\), weighted by \(\bar f(1)\) and \(\bar f(-1)\), are equal
by \eqref{eq:empirical-balance}; for (iii) insert each entry vector. (iv)
\(P(1)-P(-1)=\rho[w_F(M)-w_F(m)]+[e(1)-\rho]w_F(M)\), and \(M-m=(\pi-m)v/\bar f(1)\) by
\eqref{eq:empirical-balance}; the trading bound \(K_j\ge\rho\kappa(\pi)\DT>k\) gives full trade. (v)
\((M-\pi)-(\pi-m)=\pi(1-\pi)\lambda^2(1-2\pi)/[\bar f(1)\bar f(-1)]\). (vi) Average \(w_j=pF(p)+\int_p^\infty[1-G_j][1-F]\,dy\)
over the class and differentiate; each \(g_j\) is nonincreasing, so \(\chi_r\) is; at \(r_C\), \(e(1)\) drops from \(1\)
to \(\rho\).
\end{derivation}}

\emph{In words.} A positive weighted difference says that the price reaches the challenger and a purchase recruits
it. Below the prior \(1/2\) a rare purchase moves the belief more, so the plain \(I_+>I_-\) can be a false positive.

\section{What order counts and prices identify (Proposition A.11, Corollary A.3)}\label{sec:empirical-identify}

{\small
\begin{paperresult}{Proposition A.11 (What order counts and prices identify)}
Let full two-sided trade hold.
\begin{enumerate}[label=(\alph*),leftmargin=2em,nosep]
\item \emph{A known prior other than \(1/2\).} If \(\pi\ne1/2\) is known, the counts identify the slot:
\(
\lambda=\frac{\E[X]}{2\pi-1},\ b=\frac{v-\lambda}{2},\ \beta=\frac{v-\lambda}{2(1-\lambda)},\ \vartheta=\frac{\lambda}{v}.
\)
Over independent order opportunities, the mean order divided by \(2\pi-1\) estimates \(\lambda\) without bias, with variance \((v-\E[X]^2)/(2\pi-1)^2\) divided by the number of opportunities. If also \(\pi<\tau\), the sale impact identifies \(\rho\DT=\bar f(-1)\,I_-/[\pi(1-\pi)\lambda]\).
\item \emph{The prior \(1/2\).} At the prior \(1/2\) the order law is \((v/2,1-v,v/2)\), so the counts identify \(v\) and no other function of the slot. A share \(\vartheta'\in(0,1)\) belongs to exactly one pair \((\lambda',\beta')\) with the volume \(v\), namely \(\lambda'=v\vartheta'\) and \(\beta'=v(1-\vartheta')/[2(1-v\vartheta')]\). This pair is a slot of the model, \(\beta'\le1/3\), exactly when \(\vartheta'\ge3-2/v\). Let \(\tau>1/2\). Under full two-sided trade at this slot, entry is \(\rho+(1-\rho)(v/2)\,\1\{(1+\vartheta')/2\ge\tau\}\) when an indifferent high-cost challenger prepares. Full two-sided trade at this slot is an equilibrium exactly when \(\vartheta'\le1-2k/(\rho\DT)\) if an indifferent investor may mix over its best orders, and exactly when \(\vartheta'<1-2k/(\rho\DT)\) if an indifferent investor abstains.
\item \emph{The sale price.} If entry after a sale and after no order is \(\rho\) and \(\rho\DT\) is known, the sale impact gives \(\pi(1-\pi)\lambda=\bar f(-1)\,I_-/(\rho\DT)\). With the counts,
\(
\lambda=\frac{2\bar f(-1)\,I_-}{\rho\DT}+\sqrt{\bigl(\frac{2\bar f(-1)\,I_-}{\rho\DT}\bigr)^2+\E[X]^2},\
\pi=\frac12\bigl(1+\frac{\E[X]}{\lambda}\bigr),
\)
and then \(b\) and \(\beta\) as in part (a). At the prior \(1/2\), \(\lambda=4\bar f(-1)\,I_-/(\rho\DT)\) and \(\vartheta=2I_-/(\rho\DT)\).
\end{enumerate}
\end{paperresult}}

{\small
\begin{paperresult}{Corollary A.3 (Counts under purchases only)}
Take an equilibrium in which class \(L\) abstains and class \(H\) buys with probability \(\sigma_H(1)>0\). Write \(s_+=\lambda\sigma_H(1)\).
\begin{enumerate}[label=(\alph*),leftmargin=2em,nosep]
\item Then \(\bar f(-1)=b\), \(\bar f(1)=b+\pi s_+\), \(\E[X]=\pi s_+\) and \(v=2b+\pi s_+\). At every prior, the counts identify \(b=(v-\E[X])/2\) and the share \(\E[X]/v\) of nonzero orders that the informed investor places. At a known prior they identify \(s_+=\E[X]/\pi\). If class \(H\) buys with probability one, \(\lambda=\E[X]/\pi\).
\item Suppose that an indifferent investor may mix over its best orders and that class \(H\) mixes, \(0<\sigma_H(1)<1\). For every \(\lambda'\in[s_+,1-3b]\), the slot \((\lambda',\beta')\) with \(\beta'=b/(1-\lambda')\), with \(\sigma'_H(1)=s_+/\lambda'\) and with class \(L\) abstaining, has an equilibrium with the same class-conditional order laws, and so the same posteriors, prices, entry and acquisition probabilities. The counts therefore identify \(b\) and, at a known prior, \(s_+\). Even with the prices, they do not identify \(\lambda\). On a mixed purchase, the indifference of class \(H\) gives \(s_+=(b/\pi)\bigl[e(1)(1-\pi)\DT/k-1\bigr]\), which is \((b/\pi)\bigl[(1-\pi)\DT/k-1\bigr]\) when every challenger prepares after a purchase.
\end{enumerate}
\end{paperresult}}

{\small\begin{derivation}[proof]
A.11: \(v=2b+\lambda\) and \(\E[X]=\lambda(2\pi-1)\); \(X^2=\1\{X\ne0\}\) gives \(\operatorname{Var}X=v-\E[X]^2\). At
\(\pi=1/2\), \(\bar f(\pm1)=v/2\), \(M'=(1+\vartheta')/2\), and full trade needs \(\rho\DT m'\ge k\). With
\(q=\bar f(-1)I_-/(\rho\DT)=\pi(1-\pi)\lambda\), \(\E[X]^2=\lambda^2-4q\lambda\), whose positive root is the display.
A.3: the class laws \(f_H=(b,1-2b-s_+,b+s_+)\), \(f_L=(b,1-2b,b)\) depend only on \(b\) and \(s_+\); the indifference
\(e(1)\DT(1-\pi)b=k(b+\pi s_+)\) gives \(s_+\).
\end{derivation}}

\section{Orders predict the winner (Proposition A.12)}\label{sec:empirical-winner}

Notation: \(W_j(F)=\Pr(R<\theta)\), \(\theta\sim G_j\); \(\Omega=W_H-W_L\); \(W_F(\mu)=W_L+\mu\Omega\); \(\Gamma(x)=\int_p^xD\,dy\).

{\small
\begin{paperresult}{Proposition A.12 (Orders predict the winner)}
\begin{enumerate}[label=(\roman*),leftmargin=2em,nosep]
\item In every equilibrium, the probability that the challenger acquires the target after the order \(x\), given preparation, is \(W_F(\mu(x))\), and \(\E[\1_{\rm acq}\mid X=x]=e(x)\,W_F(\mu(x))\). Also \(\Omega(F)=\E_F[D(R)]\ge0\), and \(\Omega(F)>0\) exactly when \(\Pr_F\{D(R)>0\}>0\).
\item Under full two-sided trade, \(W_F(M)-W_F(m)=(M-m)\,\Omega(F)\ge0\), positive exactly when \(\Pr_F\{D(R)>0\}>0\). The prices rise strictly with the order, so the win probability of a prepared challenger rises with the price, strictly when \(\Pr_F\{D(R)>0\}>0\).
\item Under full two-sided trade with \(\pi<\tau\),
\(\Cov(X,\1_{\rm prep})=[e(1)-\rho]\,\bar f(1)\,(1-\E[X])\),
\(\Cov(X,\1_{\rm acq})=2\pi(1-\pi)\rho\lambda\,\Omega(F)+[e(1)-\rho]\,\bar f(1)\,(1-\E[X])\,W_F(M)\).
If \(B_F(M)<c_H\), entry is \(\rho\) after every order, \(\Cov(X,\1_{\rm prep})=0\), and \(\Cov(X,\1_{\rm acq})=2\pi(1-\pi)\rho\lambda\,\Omega(F)\), which is positive exactly when \(\Pr_F\{D(R)>0\}>0\).
\item Under full two-sided trade with \(\pi<\tau\),
\(\E[\1_{\rm acq}\mid X=1]-2\,\E[\1_{\rm acq}\mid X=0]+\E[\1_{\rm acq}\mid X=-1]=\rho\,(M+m-2\pi)\,\Omega(F)+[e(1)-\rho]\,W_F(M)\).
At the prior \(1/2\), \(M+m=1=2\pi\). If also an indifferent high-cost challenger prepares, then \(\Cov(X,\1_{\rm acq})=\rho\lambda\,\Omega(F)/2+(1-\rho)\,\chi_F\,\bar f(1)\,W_F(M)\), and the second difference equals \((1-\rho)\,\chi_F\,W_F(M)\), which is positive exactly when a purchase recruits.
\item Let \(F_r\) be a path, hold the prior and the slot fixed, and let class \(H\) always outbid the incumbent: \(W_H(F_r)=1\) at every strength. At strengths with full two-sided trade, the winner gap \(W_{F_r}(M)-W_{F_r}(m)=(M-m)\,[1-W_L(F_r)]\) is nondecreasing in \(r\), and it is larger at \(r'\) than at \(r<r'\) exactly when \(W_L(F_{r'})<W_L(F_r)\).
\item If \(F_r\) is uniform on \([0,r]\) with \(r\ge p\), then \(\Omega(F_r)=\Gamma(r)/r\), and its right derivative in \(r\) is \([D(r)-\Omega(F_r)]/r\). So \(\Omega(F_r)\), and with it the full-trade winner gap \((M-m)\,\Omega(F_r)\) at a fixed prior and slot, is nondecreasing on an interval of strengths exactly when \(D(r)\ge\Omega(F_r)\) at almost every strength of the interval. It rises strictly on an interval where \(D(r)>\Omega(F_r)\) at almost every strength.
\end{enumerate}
\end{paperresult}}

{\small\begin{derivation}[proof]
Preparation is independent of \((\theta,R)\) given the class, so a prepared challenger wins with \(W_F(\mu(x))\), and
\(\Omega(F)=\int(G_L-G_H)\,dF=\E_F[D(R)]\). For (iii), \(\Cov(X,Y)=\Cov(X,\E[Y\mid X])\),
\(\Cov(X,c\1\{X=1\})=c\bar f(1)(1-\E[X])\), and \(\Cov(X,\mu(X))=\pi\lambda-\pi\lambda(2\pi-1)=2\pi(1-\pi)\lambda\).
(iv) \(W_F\) is affine, so its second difference is \((M+m-2\pi)\Omega\). (vi) For the uniform law,
\(\Omega(F_r)=r^{-1}\int_0^rD\,dy=\Gamma(r)/r\); apply the quotient rule.
\end{derivation}}

\section{Announcement returns (Proposition A.13)}\label{sec:empirical-announce}

{\small
\begin{paperresult}{Proposition A.13 (Announcement returns)}
Take any equilibrium, at any prior, under either challenger rule.
\begin{enumerate}[label=(\roman*),leftmargin=2em,nosep]
\item After the order \(x\), the price changes by \(\Delta^1_x=[1-e(x)]\,w_F(\mu(x))\) when the event reports preparation, and by \(\Delta^0_x=-e(x)\,w_F(\mu(x))\) when it reports no preparation. The second report has positive probability exactly when \(e(x)<1\).
\item After an order \(x\) with \(e(x)<1\), the mean change over both reports, with their probabilities, is zero.
\item Pool all orders, and take the sample of preparation reports. Its mean price change is
\(
\sum_x\Pr(X=x\mid\text{preparation})\,\Delta^1_x=\frac{\sum_x\bar f(x)\,e(x)\,[1-e(x)]\,w_F(\mu(x))}{\sum_x\bar f(x)\,e(x)},
\)
which is positive exactly when entry is below one after some order, that is, exactly when preparation is uncertain. If \(B_F(\pi)<c_H\), entry is \(\rho\) after some order, and the mean is positive.
\item The expected price before the order opportunity minus the price after no order is the weighted impact difference, \(\E[P(X)]-P(0)=\bar f(1)I_+-\bar f(-1)I_-\), which parts (i) to (iii) of Proposition A.10 evaluate. Under full two-sided trade the posterior after no order is the prior.
\end{enumerate}
\end{paperresult}}

{\small\begin{derivation}[proof]
After a report of preparation the price is \(t_0+w_F(\mu(x))\); after no preparation it is \(t_0\). Subtract
\(P(x)=t_0+e(x)w_F(\mu(x))\); the two changes average to zero with weights \(e(x)\), \(1-e(x)\). Pool with Bayes'
rule. If \(B_F(\pi)<c_H\), then \(\sum_x\bar f(x)B_F(\mu(x))=B_F(\pi)<c_H\), so some order has \(\mu(x)<\tau\).
\end{derivation}}

\begin{watchout}
With a hidden price (simple deterrence) the weighted difference, the entry component and \(\Cov(X,\1_{\rm prep})\) are
zero, so these separate the mechanism; the winner and announcement predictions hold in any informative market. Each
strength is a separate economy: hold the state fixed.
\end{watchout}

\begin{chaptersummary}
Under full two-sided trade, the weighted impact difference is positive exactly when a purchase recruits; the balance
identity removes the class term. Counts identify the informed presence at a known prior other than \(1/2\), and only
the volume at \(1/2\). The order predicts the winner through the class, and the mean return on preparation reports is
positive whenever preparation is uncertain.
\end{chaptersummary}
